%==============================================
% VOLUME I: Zero to Bit
%==============================================
% What this volume covers:
%   The foundational question — how is information encoded in digital
%   systems at all, starting from a voltage difference?
%
%   Path: voltage → binary → numbers → characters → bytes → pointers → integrity → serialization → compression
%
% Assumes: nothing beyond basic programming familiarity
%==============================================

\begin{volumeintro}
\lettrine[lines=3]{B}{efore we} can understand how arrays store elements, we must first understand how computers represent information itself. Every piece of data---numbers, text, images, instructions---exists as patterns of bits, and every bit begins its life as a voltage. This volume follows that single thread: from a voltage difference across a transistor, through binary encoding, to the numbers, characters, and bit-level structures that everything else in this work is built from.

\begin{quote}
\textit{``The choice of representation is often more important than the choice of algorithm.''}

\hfill--- \textsc{Donald E. Knuth}
\end{quote}
\end{volumeintro}

% ──────────────────────────────────────────────────────────────────────
% STAGE 1 — What Data Is and What It Means to Represent It
%
%   Before a single bit appears, we need to understand what data is,
%   and what it means for one physical thing to stand for another.
%   These two chapters are the philosophical foundation everything else
%   rests on.
% ──────────────────────────────────────────────────────────────────────

\chapter{The Nature of Data}
\label{ch:nature-of-data}

What data actually is before it is encoded: the distinction between data,
information, and meaning; raw measurements versus interpreted values; data
as the record of a distinction or a difference; the role of an observer or
interpreter in turning a physical trace into data; examples spanning natural
``data'' (tree rings, sediment layers, starlight) and human-made data (tally
marks, ledgers, sensor readings); why data alone is inert without a scheme
to read it; the seed of the idea that everything downstream---bits, numbers,
arrays---is data wrapped in successive layers of agreed-upon meaning.
